\documentclass[11pt]{article}
\usepackage[a4paper,margin=1in]{geometry}
\usepackage{amsmath,amssymb,mathtools,bm}
\usepackage{enumitem,booktabs,array}
\usepackage{tcolorbox}
\usepackage{hyperref}
\usepackage[protrusion=true,expansion=false]{microtype}
\hypersetup{colorlinks=true,linkcolor=blue,urlcolor=blue,citecolor=blue}
\setlist[itemize]{topsep=2pt,itemsep=2pt}
\setlist[enumerate]{topsep=2pt,itemsep=2pt}
\newtcolorbox{keybox}{colback=blue!3!white,colframe=blue!40!black,boxrule=0.4pt,arc=1mm}
\newtcolorbox{alertbox}{colback=red!3!white,colframe=red!40!black,boxrule=0.4pt,arc=1mm}
\newtcolorbox{answerbox}{colback=green!3!white,colframe=green!40!black,boxrule=0.4pt,arc=1mm}
\newtcolorbox{drillbox}{colback=yellow!5!white,colframe=orange!60!black,boxrule=0.4pt,arc=1mm}
\newcommand{\EV}{\mathbb{E}}
\newcommand{\Var}{\mathrm{Var}}
\newcommand{\Cov}{\mathrm{Cov}}
\newcommand{\blank}[1]{\underline{\hspace{#1}}}

\title{E10-02: Lerner \& Malmendier (2013, RFS) \\
\large ``With a Little Help from My (Random) Friends: Success and Failure in Post-Business School Entrepreneurship'' --- Active Recall Workbook}
\author{Empirical Corporate Finance II --- Exam Prep}
\date{\today}
\begin{document}\maketitle

\begin{keybox}
\textbf{Paper in one sentence.} Exploiting the random assignment of Harvard Business School students to sections and the ensuing cohort peer groups, LM show that having more peers with prior entrepreneurial experience \emph{reduces} the probability that a student becomes an entrepreneur but \emph{improves} the average success conditional on entering, evidencing a peer-learning channel in career choice.

\textbf{Placement.} Canonical quasi-experimental paper on peer effects in entrepreneurship; highlights the role of information and social learning in professional networks.
\end{keybox}

\section*{Round 1: Complete Walk-Through}

\subsection*{1.1 Setting}
HBS assigns incoming MBAs to sections of $\sim$90 students each, quasi-randomly conditional on gender and nationality. Students share classes, projects, and social networks within the section. Section assignment provides exogenous variation in peer composition.

\subsection*{1.2 Data}
\begin{itemize}
\item HBS MBA classes 1997--2004.
\item Detailed pre-MBA backgrounds: prior industry, entrepreneurial experience, financial industry.
\item Post-MBA career choices and outcomes.
\end{itemize}

\subsection*{1.3 Empirical design}
\begin{itemize}
\item Outcome: whether student starts a firm within 3 years of graduation.
\item Treatment: fraction of section peers with prior entrepreneurial experience.
\item Regressions include year-cohort fixed effects; identification from within-cohort section-level variation.
\end{itemize}

\subsection*{1.4 Headline results}
\begin{itemize}
\item More entrepreneurial peers reduces probability of entering entrepreneurship by $\sim$1--2 pp per peer (vs.\ baseline $\sim$10\%).
\item But conditional on entering, performance (firm survival, growth) improves.
\item Interpretation: peer information corrects overoptimism; those who still enter are better-matched.
\item Effect is asymmetric: peers with financial-industry background increase entry but do not affect performance.
\end{itemize}

\subsection*{1.5 Mechanism}
\begin{itemize}
\item Peer learning and screening: experienced peers share honest information about entrepreneurial costs.
\item Reduces excessive entry among over-optimistic students.
\item Raises quality of those who do enter.
\end{itemize}

\subsection*{1.6 Implications}
\begin{itemize}
\item Peer effects in career choice are significant and information-based.
\item Entrepreneurship support programs may benefit from experienced mentors.
\item Challenges simple ``more exposure = more entrepreneurs'' narratives.
\end{itemize}

\subsection*{1.7 Caveats}
\begin{alertbox}
\begin{itemize}
\item Harvard MBA students are non-representative.
\item Sample size within-cohort limits power.
\item Outcomes (firm survival) are crude proxies for success.
\end{itemize}
\end{alertbox}

\section*{Round 2: Level 1 Fill-in}
\begin{enumerate}
\item Exogenous variation: HBS \blank{2cm} assignment.
\item Sample: MBA classes \blank{1cm}--\blank{1cm}.
\item Entry effect: $-$\blank{1cm}--\blank{1cm} pp per peer.
\item Performance effect: \blank{2cm} improves.
\item Channel: peer \blank{2cm} corrects overoptimism.
\end{enumerate}

\section*{Round 3: Level 2 Prompts}
\begin{enumerate}
\item Describe the HBS section setting and identification.
\item Report main findings on entry and performance.
\item Discuss the peer-learning mechanism.
\item Differentiate from simple ``exposure increases entry'' views.
\item Evaluate external validity.
\end{enumerate}

\section*{Round 4: Minimal Scaffolding}
From a blank page: (i) setting, (ii) data, (iii) findings, (iv) mechanism, (v) implications.

\section*{Round 5: Blank Page}
Essay: peer effects in entrepreneurial career choice.

\vspace{6pt}
\begin{drillbox}
\textbf{Speed Drill.} (1) Identification. (2) Sample. (3) Entry effect. (4) Performance effect. (5) Mechanism.
\end{drillbox}

\section*{Quick Reference \& Answer Key}
\begin{answerbox}
\begin{itemize}
\item \textbf{Setting.} Randomized HBS sections.
\item \textbf{Entry.} Falls with entrepreneurial peers.
\item \textbf{Performance.} Rises conditional on entry.
\item \textbf{Lesson.} Peer learning corrects overoptimism.
\end{itemize}
\end{answerbox}

\begin{keybox}
\textbf{30-second exam pitch.} Lerner \& Malmendier (2013) use the quasi-random assignment of HBS MBA students to sections (1997--2004) to identify peer effects on entrepreneurial career choice. With within-cohort, cross-section variation in peers with prior entrepreneurial experience, they find a counter-intuitive result: more entrepreneurial peers reduce the probability that a given student enters entrepreneurship by 1--2 pp per peer (vs.\ a 10\% baseline), yet conditional on entering, performance (survival, growth) improves. The interpretation is a peer-learning channel: experienced peers honestly share costs and risks, deterring over-optimistic entry while raising the quality of those who do enter. The paper is a clean quasi-experimental demonstration of information-based peer effects in career choice and has influenced entrepreneurship-mentorship and business-school-program design debates.
\end{keybox}

\end{document}
